\documentclass[aspectratio=169]{beamer}
\usetheme{Madrid}
\usecolortheme{whale}
\usepackage[utf8]{inputenc}
\usepackage[spanish]{babel}
\usepackage{amsmath,amsfonts,amssymb}
\usepackage{graphicx}
\usepackage{booktabs}
\usepackage{xcolor}

\title{Métodos de Análisis de Datos 1}
\subtitle{Clase 11: Análisis Bivariado}
\author{Depto.\ de Matemática -- UNS}
\institute{Universidad Nacional del Sur}
\date{1er Cuatrimestre 2026}

\AtBeginSection[]{
  \begin{frame}{Contenidos}
    \tableofcontents[currentsection]
  \end{frame}
}

\begin{document}

\begin{frame}
  \titlepage
\end{frame}

\begin{frame}{Contenidos}
  \tableofcontents
\end{frame}

%=============================================================
\section{Introducción al análisis bivariado}
%=============================================================

\begin{frame}{¿Qué es el análisis bivariado?}
  \begin{block}{Objetivo}
    Estudiar la relación o asociación entre \textbf{dos variables} medidas sobre las mismas unidades.
  \end{block}
  \bigskip
  \begin{center}
    \begin{tabular}{lll}
      \toprule
      \textbf{Variable 1} & \textbf{Variable 2} & \textbf{Herramienta} \\
      \midrule
      Cuantitativa & Cuantitativa & Scatter plot, correlación, regresión \\
      Cualitativa & Cualitativa & Tabla de contingencia, $\chi^2$ \\
      Cualitativa & Cuantitativa & Boxplot por grupos, ANOVA \\
      \bottomrule
    \end{tabular}
  \end{center}
  \bigskip
  \begin{alertblock}{Cuidado}
    \textbf{Asociación no implica causalidad}. Siempre interpretar en contexto.
  \end{alertblock}
\end{frame}

%=============================================================
\section{Variables cuantitativas: correlación}
%=============================================================

\begin{frame}{Gráfico de dispersión (scatter plot)}
  \begin{itemize}
    \item Representación de pares $(x_i, y_i)$ en el plano cartesiano.
    \item Permite detectar:
      \begin{itemize}
        \item Dirección de la relación: positiva, negativa, ninguna.
        \item Forma: lineal, curvilínea.
        \item Fuerza: dispersión alrededor de la tendencia.
        \item Outliers bivariados: puntos aislados del patrón general.
      \end{itemize}
    \item Se puede agregar una \textbf{línea de tendencia} (regresión simple) para visualizar la relación lineal.
  \end{itemize}
\end{frame}

\begin{frame}{Correlación de Pearson}
  \begin{block}{Coeficiente de correlación de Pearson ($r$)}
    \[
      r = \frac{\sum_{i=1}^n (x_i - \bar{x})(y_i - \bar{y})}
               {\sqrt{\sum_{i=1}^n (x_i - \bar{x})^2 \cdot \sum_{i=1}^n (y_i - \bar{y})^2}}
    \]
  \end{block}
  \begin{itemize}
    \item $-1 \le r \le 1$.
    \item $r = 1$: correlación lineal perfecta positiva.
    \item $r = -1$: correlación lineal perfecta negativa.
    \item $r = 0$: sin correlación lineal (puede haber relación no lineal).
    \item \textbf{Supuestos}: normalidad bivariada, relación lineal, sin outliers influyentes.
  \end{itemize}
\end{frame}

\begin{frame}{Interpretación del coeficiente $r$}
  \begin{center}
    \begin{tabular}{cc}
      \toprule
      \textbf{Valor de $|r|$} & \textbf{Interpretación} \\
      \midrule
      $0.00 - 0.19$ & Correlación muy débil \\
      $0.20 - 0.39$ & Correlación débil \\
      $0.40 - 0.59$ & Correlación moderada \\
      $0.60 - 0.79$ & Correlación fuerte \\
      $0.80 - 1.00$ & Correlación muy fuerte \\
      \bottomrule
    \end{tabular}
  \end{center}
  \bigskip
  \begin{alertblock}{Significación estadística}
    Con $n$ grande, correlaciones pequeñas pueden ser estadísticamente significativas pero no prácticamante relevantes. Siempre reportar $r$ y su $p$-valor juntos.
  \end{alertblock}
\end{frame}

\begin{frame}{Correlación de Spearman}
  \begin{block}{Coeficiente $\rho$ de Spearman}
    \[
      \rho = 1 - \frac{6 \sum_{i=1}^n d_i^2}{n(n^2-1)}
    \]
    donde $d_i = \text{rango}(x_i) - \text{rango}(y_i)$.
  \end{block}
  \bigskip
  \textbf{Ventajas sobre Pearson:}
  \begin{itemize}
    \item No asume normalidad ni linealidad: mide asociación monótona.
    \item \textbf{Robusta} frente a outliers.
    \item Aplicable a variables ordinales.
  \end{itemize}
  \bigskip
  \begin{exampleblock}{Recomendación}
    Usar Spearman cuando los supuestos de Pearson no se cumplen o con variables ordinales.
  \end{exampleblock}
\end{frame}

%=============================================================
\section{Variables cualitativas: tablas de contingencia}
%=============================================================

\begin{frame}{Tabla de contingencia}
  Para dos variables categóricas $A$ (con $r$ categorías) y $B$ (con $c$ categorías):
  \[
    \begin{array}{c|cccc|c}
      & B_1 & B_2 & \cdots & B_c & \text{Total} \\
      \hline
      A_1 & n_{11} & n_{12} & \cdots & n_{1c} & n_{1\cdot} \\
      A_2 & n_{21} & n_{22} & \cdots & n_{2c} & n_{2\cdot} \\
      \vdots & & & & & \vdots \\
      A_r & n_{r1} & n_{r2} & \cdots & n_{rc} & n_{r\cdot} \\
      \hline
      \text{Total} & n_{\cdot 1} & n_{\cdot 2} & \cdots & n_{\cdot c} & n \\
    \end{array}
  \]
  \begin{itemize}
    \item $n_{ij}$: frecuencia observada en celda $(i,j)$.
    \item Frecuencias \textbf{marginales}: $n_{i\cdot}$ (fila) y $n_{\cdot j}$ (columna).
  \end{itemize}
\end{frame}

\begin{frame}{Test de independencia $\chi^2$}
  \begin{block}{Hipótesis}
    $H_0$: las variables $A$ y $B$ son independientes. \quad $H_1$: hay asociación.
  \end{block}
  \begin{block}{Estadístico}
    \[
      \chi^2 = \sum_{i,j} \frac{(O_{ij} - E_{ij})^2}{E_{ij}}, \quad
      E_{ij} = \frac{n_{i\cdot} \cdot n_{\cdot j}}{n}
    \]
    Bajo $H_0$: $\chi^2 \sim \chi^2_{(r-1)(c-1)}$.
  \end{block}
  \begin{alertblock}{Condición de validez}
    Las frecuencias esperadas $E_{ij} \ge 5$ en al menos el 80\% de las celdas. Si no se cumple: combinar categorías o usar test exacto de Fisher.
  \end{alertblock}
\end{frame}

\begin{frame}{Medidas de asociación para variables categóricas}
  \begin{center}
    \begin{tabular}{lll}
      \toprule
      \textbf{Medida} & \textbf{Aplicación} & \textbf{Rango} \\
      \midrule
      V de Cramér & Nominal $\times$ nominal & $[0, 1]$ \\
      Phi ($\phi$) & Tablas $2\times2$ & $[0, 1]$ \\
      Gamma de Goodman & Ordinal $\times$ ordinal & $[-1, 1]$ \\
      $\tau_b$ de Kendall & Ordinal $\times$ ordinal & $[-1, 1]$ \\
      OR (Odds Ratio) & Tablas $2\times 2$ & $(0, +\infty)$ \\
      \bottomrule
    \end{tabular}
  \end{center}
  \bigskip
  \begin{exampleblock}{V de Cramér}
    $V = \sqrt{\chi^2 / (n \cdot \min(r-1, c-1))}$. Interpretación: 0 = sin asociación, 1 = asociación perfecta.
  \end{exampleblock}
\end{frame}

%=============================================================
\section{Cuantitativa vs.\ cualitativa}
%=============================================================

\begin{frame}{Comparación de grupos}
  \textbf{Variable cuantitativa por grupos (variable cualitativa):}
  \bigskip
  \begin{itemize}
    \item \textbf{Boxplot por grupos}: comparar distribución de $Y$ para cada nivel de $X$.
    \item \textbf{Estadísticos por grupo}: media, mediana, DE para cada categoría.
    \item \textbf{Tests estadísticos}:
      \begin{itemize}
        \item 2 grupos, normalidad: test $t$ de Student.
        \item 2 grupos, no normal: Wilcoxon-Mann-Whitney.
        \item $k$ grupos, normalidad: ANOVA (se verá en U3).
        \item $k$ grupos, no normal: Kruskal-Wallis.
      \end{itemize}
  \end{itemize}
  \bigskip
  \begin{alertblock}{Siempre verificar primero}
    Normalidad dentro de cada grupo (Shapiro-Wilk) y homocedasticidad (test de Levene o Bartlett) antes de usar tests paramétricos.
  \end{alertblock}
\end{frame}

%=============================================================
\section{Resumen}
%=============================================================

\begin{frame}{Resumen de la clase}
  \begin{block}{Análisis bivariado}
    \begin{itemize}
      \item Cuant.-Cuant.: scatter plot, Pearson, Spearman.
      \item Cual.-Cual.: tabla de contingencia, $\chi^2$, V de Cramér.
      \item Cual.-Cuant.: boxplot por grupos, tests de comparación.
    \end{itemize}
  \end{block}
  \bigskip
  \begin{block}{Próximas clases}
    Clase 12: Visualización con Python — \texttt{matplotlib}, \texttt{seaborn}, buenas prácticas. \\
    Clase 13: Visualización con R — \texttt{ggplot2} y el sistema de capas.
  \end{block}
\end{frame}

\end{document}
